\section{Source boundaries and the stationary geometric extraction}
\label{sec:source-boundary-calculus}
We give the coordinate and boundary arguments for the geometric
extraction in one place. We then establish the corresponding positive
extraction under the stationary suspension law. The latter is not the
section law: its entrance marginal has an unbounded roof bias.

\subsection{The differential and the action in invariant coordinates}
\begin{lemma}[Collision differential and action]
\label{lem:self-contained-collision-differential}
On a regular flight from $Rn_\alpha$ to $d+Rn_{\alpha'}$, set
$c=\sqrt{1-p^2}$, $c'=\sqrt{1-(p')^2}$ and $v_0=\tau_R/R$.
In the coordinates $(\alpha,p)$ the derivative is
\[
 DT_R=-\begin{pmatrix}
 (v_0+c)/c'&v_0/(cc')\\
 v_0+c+c'&(v_0+c')/c
 \end{pmatrix},\qquad \det DT_R=1,
\]
and
\begin{equation}\label{eq:self-contained-action-differential}
 d\tau_R=Rp'\,d\alpha'-Rp\,d\alpha.
\end{equation}
These formulas use the outgoing value of $p'$ after reflection.
\end{lemma}
\begin{proof}
Write $\mathcal J$ for counterclockwise rotation by $\pi/2$, and
$V=cn_\alpha+pt_\alpha$ for the unit flight velocity. If $\beta$ is
its direction angle on a local chart, then
\[
 d\beta=d\alpha+dp/c,\qquad dV=\mathcal JV\,d\beta.
\]
At arrival, $V=-c'n_{\alpha'}+p't_{\alpha'}$: reflection preserves
the tangential component and reverses the normal component. Thus
$d\beta=d\alpha'-dp'/c'$. Differentiating the flight equation and
taking its scalar product with $\mathcal JV$ gives
\[
 -Rc'\,d\alpha'=Rc\,d\alpha+\tau_R\,d\beta,
\]
since $\mathcal JV\cdot t_\alpha=c$ and
$\mathcal JV\cdot t_{\alpha'}=-c'$. Substitution gives the first
row of the displayed matrix. The identity
$dp'=c'(d\alpha'-d\beta)$ gives its second row. Its determinant is
\[
 \frac{(v_0+c)(v_0+c')-v_0(v_0+c+c')}{cc'}=1.
\]
Taking instead the scalar product of the differentiated flight equation
with $V$ gives \eqref{eq:self-contained-action-differential}.
The lattice center is constant on this branch. Summing the action
formula gives the two derivatives of $L_m$ used in
Lemma~\ref{lem:roof-transverse-direction}. This derivation also
specifies the sign and coordinate conversion in the fixed-table
specialization of \cite[Section 3.1]{SYZ}.
\end{proof}

\subsection{A complete first-hit guard}
Let $\tau_{\max}$ be a uniform horizon from the earlier geometric
estimates, and choose once and for all
\begin{equation}\label{eq:explicit-finite-center-set}
 \mathcal D=\{d\in\Lambda\setminus\{0\}:
                        |d|\le\tau_{\max}+2R_++1\}.
\end{equation}
Here $\Lambda$ is the physical triangular lattice, not a rescaled
collision coordinate lattice.

\begin{lemma}[First-hit transitions and uniform chart coefficients]
\label{lem:complete-first-hit-guard}
The set \eqref{eq:explicit-finite-center-set} is finite and contains
every next collision center for all $R\in I$. Away from initial grazing
and the zero sets of all $\Delta_{d,R}$, the first-hit label is locally
constant and the selected positive entry root is smooth. In four bounded
angular half-angle charts, the numerator of each discriminant has degree
at most eight, its positive denominator lies in $[1,16]$, and the maximum
modulus of its coefficients is bounded below uniformly in $R$, the
chart and $d\in\mathcal D$.
\end{lemma}
\begin{proof}
A flight of length at most $\tau_{\max}$ from one boundary point to the
next gives $|d|\le\tau_{\max}+2R_+$, proving the first assertion.
For a candidate let $a=d-Rn_\alpha$, $A=a\cdot V$ and
$\Delta=R^2-|a|^2+A^2$. Since $|d|\ge1$ and $R\le47/100$,
\[
 |a|\ge1-R>R,\qquad A^2-\Delta=|a|^2-R^2>0.
\]
On $\Delta>0$ the entry root is $A-\sqrt\Delta$. It is positive
exactly when $A>0$. Its sign cannot change while $\Delta>0$, and a
positive root cannot pass through zero. Two distinct positive entry
roots cannot be equal: the corresponding point on the ray would lie
on two disjoint closed disks. Consequently the ordering of all admissible
entry roots persists locally until a candidate discriminant vanishes.
The margin in \eqref{eq:explicit-finite-center-set} prevents a center
from entering through the edge of the candidate list. The outgoing
half-ray cannot reenter the initial convex disk. This accounts for all
first-hit transitions, including the side of a tangency on which the
candidate is no longer hit.

Take chart centers $\alpha_j=j\pi/2$, $0\le j\le3$, and use
$\alpha=\alpha_j+2\arctan t$ on $|t|<1$; their open images cover the
angular circle. Bounds may be taken on the closed chart intervals.
Let $s=\tan(\varphi/2)\in[-1,1]$, where $p=\sin\varphi$.
Then
\[
 p=\frac{2s}{1+s^2},\qquad c=\frac{1-s^2}{1+s^2},\qquad
 a\mathbin{\times}V
  =c(d\mathbin{\times}n_\alpha)
                  +p(d\mathbin{\times}t_\alpha-R).
\]
The last expression has denominator $(1+t^2)(1+s^2)$ and a
numerator of total degree at most four. Thus the discriminant has
denominator $(1+t^2)^2(1+s^2)^2\in[1,16]$ and a numerator of total
degree at most eight. The coordinate Jacobian of $\nu$ is bounded
above by $1/\pi$, since $|d\alpha/dt|\le2$ and $|dp/ds|\le2$.

For each fixed $d,j,R$, the numerator is not identically zero. Fix an
initial point in the chart. Its vector $a$ has length greater than $R$;
the square of its cross product with directions in an open angular
interval cannot be the constant $R^2$. The maximum modulus of the
finitely many polynomial coefficients is a continuous, strictly positive
function of $R$. Its minimum on $I$ is positive. Taking the minimum
over the finite candidate and chart sets proves the uniform coefficient
bound. This justifies the uniform use of
Lemma~\ref{lem:elementary-polynomial-sublevel}, including near tangency.
\end{proof}

\begin{lemma}[Partition and zero extension]
\label{lem:source-partition-zero-extension}
Fix $n\le L$ and $\varepsilon>0$. Up to a $\nu$-null set, the source
$Y_R^*\cap\{N_{n,R}\le L\}$ is the union of finitely many disjoint open
pieces $D$ relative to the collision cylinder. On each piece the center
itinerary, section decisions, mark collision indices, output label and
$m=N_{n,R}$ are fixed. The source functions $b_D$ in the proof of
Theorem~\ref{thm:geometric-extraction-budget}, and the vector-field
products used there, extend by zero with the regularity needed for both
integrations by parts. In particular there are no distributions
supported on itinerary or section boundaries in those identities.
\end{lemma}
\begin{proof}
At each collision record the first-hit center in the finite set
\eqref{eq:explicit-finite-center-set} and the section membership decision.
There are finitely many such records through $L$ collisions. Remove
all encountered tangencies, grazing points and section boundaries.
Their union is null: the one-step discriminant zero sets are null by
Lemma~\ref{lem:complete-first-hit-guard} and the polynomial sublevel
bound, section sides are null, and the collision map preserves $\nu$.
The regular decision sets are open. They need not be connected for the
argument; all their components with one record are treated together.
On each such set every return index and output label is fixed. No
artificial angular seam is introduced in this partition.

Here is the boundary extension explicitly. At a boundary point choose
the first singular collision or the first section decision at which the
record can change. All preceding regular branches persist on a
neighborhood. At a first-hit tangency the input discriminant of the
preceding flight is zero. At an initial grazing point the initial guard
is zero; grazing produced by a flight is already its tangency. At a
section decision the corresponding rectangle-side guard is zero.
In each case, at the earliest such input time $j$, the previous maps
are continuous and $\chi_{R,\varepsilon}\circ T_R^j$ is identically
zero on a two-sided neighborhood of the boundary point. The factors
at later, possibly undefined, collision states need not be evaluated:
the product is defined to be zero throughout that neighborhood.

At a boundary of a cutoff transition inside a regular decision set,
ordinary smoothness and the flat endpoint derivatives of $\theta$
give the same matching of derivatives. The initial factor
$\theta(p/\varepsilon)$ removes the pole of $X_m$. Hence $b_D$ is
$C^2$ after zero extension, $b_DX_m$ is $C^2$, and
$X_m\operatorname{div}(b_DX_m)$ is $C^1$. Their derivatives have the
bounds of \eqref{eq:geometric-divergence-bound}. For fixed $L$ and
$\varepsilon$ their supports are separated from the excluded boundaries;
no uniform lower bound for that separation is used. The cylinder is
periodic in $\alpha$, and the supports avoid $p=\pm1$.
The distributional integration-by-parts formulas therefore have no
boundary terms. Since the pieces have disjoint interiors, their
absolute source integrals sum over a subset of the initial cylinder,
not over multiple copies of it.
\end{proof}

\subsection{The preceding section under the equilibrium law}
Write $r=r_R^*$ and $\tau^*=\tau_R^*$. On the return suspension put
\begin{equation}\label{eq:equilibrium-source-for-inversion}
 \Omega_R=\{(y,s):y\in Y_R^*,\ 0\le s<\tau^*(y)\},\qquad
 d\mathbb P_R^{\rm eq}=\frac{c_*}{\bar\tau_R}\,d\nu_R^*(y)\,ds.
\end{equation}
Its entrance marginal is $\rho_R\nu_R^*$, where
$\rho_R=c_*\tau^*/\bar\tau_R$. This is the stationary normalization
of Lemma~\ref{lem:stationary-length-bias}. The same $y$ is used for
$J_{n,R}(y)$ and for every physical trajectory observable on $\Omega_R$.

\begin{proposition}[Positive extraction with the actual entrance bias]
\label{prop:stationary-geometric-extraction}
Let $\mu_{n,R}^{\rm eq}=(J_{n,R})_*(\rho_R\nu_R^*)$.
There is a positive measure $Q\le\mu_{n,R}^{\rm eq}$ with a
compactly supported mixed density $q$ whose roof components belong to
$W^{2,1}(\R)$, and constants independent of $R,n,L,\varepsilon$ such that
\begin{align}
 \delta_{\rm eq}:=\|\mu_{n,R}^{\rm eq}-Q\|_{\TV}
 &\le C\left(e^{(an-cL)/2}
                +(L+1)^{1/2}\varepsilon^{1/32}\right),
                     \label{eq:stationary-geometric-mass}\\
 \sum_{\ell\in\Z^3}\|\partial_t^2q(\ell,\cdot)\|_1
 &\le C\exp\{C(L+1)\log(C/\varepsilon)\}=:A_2.
                     \label{eq:stationary-geometric-A2}
\end{align}
In particular $\|Q\|_{\TV}\le1$. For the section probability the same
conclusions hold with the stronger removed-mass estimate
$C(e^{an-cL}+(L+1)\varepsilon^{1/16})$.
\end{proposition}
\begin{proof}
Define
\[
 Q=(J_{n,R})_*\bigl(\rho_R\one_{\{N_{n,R}\le L\}}
                         \Gamma_{L,R,\varepsilon}\nu_R^*\bigr).
\]
All factors are nonnegative and the last two are at most one. The
one-return exponential moment and the positive lower bound for
$\bar\tau_R$ give $\sup_R\|\rho_R\|_{L^2(\nu_R^*)}<\infty$.
Cauchy--Schwarz, the cumulative return tail and
\eqref{eq:whole-itinerary-mass} consequently give
\[
 \int\rho_R\one_{\{N_n>L\}}\,d\nu_R^*
       \le C e^{(an-cL)/2},\qquad
 \int\rho_R(1-\Gamma_{L,R,\varepsilon})\,d\nu_R^*
       \le C\bigl((L+1)\varepsilon^{1/16}\bigr)^{1/2}.
\]
These imply \eqref{eq:stationary-geometric-mass}.

On a source piece of Lemma~\ref{lem:source-partition-zero-extension},
its first return collision time $j_1$ is fixed and at most $L$.
There $\tau^*=\sum_{j<j_1}\tau_R\circ T_R^j$.
The cutoff estimates and chain rule bound this function and its first
three derivatives, on the required supports, by
$C\exp\{C(L+1)\log(C/\varepsilon)\}$. Multiplying $b_D$ by
$\rho_R$ therefore preserves the divergence budget after increasing
$C$. The earliest-decision argument gives the same zero extensions.
The two source integrations by parts and disjoint-source summation
prove \eqref{eq:stationary-geometric-A2}. The argument does not assume
that the untruncated entrance bias is BV. The section-law assertion
is the unweighted case of Theorem~\ref{thm:geometric-extraction-budget}
and its cumulative-tail extension.
\end{proof}
